\documentclass[UTF8,fontset=none,a4paper,12pt]{ctexart}
\usepackage[margin=2.5cm]{geometry}
\usepackage{booktabs,longtable,array,hyperref}
\setCJKmainfont[BoldFont=SimHei,ItalicFont=KaiTi]{SimSun}
\setCJKsansfont[BoldFont=SimHei]{SimHei}
\setCJKmonofont{SimSun}
\setCJKfamilyfont{zhhei}{SimHei}
\providecommand{\heiti}{\CJKfamily{zhhei}}
\hypersetup{hidelinks}
\setlength{\emergencystretch}{3em}
\raggedbottom
\begin{document}
\begin{center}\Large\heiti AI工具使用详情\end{center}
本文件是支撑材料底稿。根据2026年全国大学生数学建模竞赛人工智能工具使用规定，交付支撑材料时应转换为 PDF，文件名为“AI工具使用详情.pdf”。

\section*{AI工具使用声明}
本参赛队在竞赛过程中使用了AI工具，主要用于代码调试、数值计算与图表制作辅助、论文结构整理和语言润色，详细使用情况见支撑材料。

\section*{表1 所用AI工具名称和版本}
{\small\setlength{\tabcolsep}{3pt}\begin{longtable}{>{\raggedright\arraybackslash}p{0.06\linewidth}>{\raggedright\arraybackslash}p{0.2\linewidth}>{\raggedright\arraybackslash}p{0.31\linewidth}>{\raggedright\arraybackslash}p{0.16\linewidth}>{\raggedright\arraybackslash}p{0.16\linewidth}}
\toprule
序号 & 工具名称 & 版本/型号 & 开发机构/公司 & 使用日期\\
\midrule\endhead
1 & Codex task（配置模型） & 具体模型版本未在权威日志中记录，待队伍复核 & OpenAI & 2026-09-10—2026-09-13\\
\bottomrule\end{longtable}}
说明：工具与阶段记录来自 .workflow/ai-usage-log.json；日志未记录更细模型型号，因此不补写未经核实的信息。

\section*{表2 具体使用目的和环节}
{\small\setlength{\tabcolsep}{3pt}\begin{longtable}{>{\raggedright\arraybackslash}p{0.06\linewidth}>{\raggedright\arraybackslash}p{0.29\linewidth}>{\raggedright\arraybackslash}p{0.15\linewidth}>{\raggedright\arraybackslash}p{0.4\linewidth}}
\toprule
序号 & 使用环节 & 具体目的 & 使用的AI工具\\
\midrule\endhead
1 & 选题与问题定义 & 整理候选赛题、数据审计和 H1 材料 & Codex task（配置模型）\\
2 & 文献与模型设计 & 形成可追溯文献证据、候选模型和验证计划 & Codex task（配置模型）\\
3 & COMPUTE & 在冻结路线下实现、调试和核验数值程序 & Codex task（配置模型）\\
4 & FIGURE & 按冻结数据生成图表、更新图注、进行 DPI/矢量/灰度检查 & Codex task（配置模型）\\
5 & PAPER & 整理已批准证据、撰写和编译论文草稿 & Codex task（配置模型）\\
\bottomrule\end{longtable}}
\section*{表3 主要提示方式与使用过程说明}
{\small\setlength{\tabcolsep}{3pt}\begin{longtable}{>{\raggedright\arraybackslash}p{0.09\linewidth}>{\raggedright\arraybackslash}p{0.32\linewidth}>{\raggedright\arraybackslash}p{0.51\linewidth}}
\toprule
序号 & 主要提示方式 & 使用过程说明\\
\midrule\endhead
1 & 以阶段 Skill、输入文件、唯一写入目录和停止条件为约束的任务提示 & 先读取阶段交接与冻结决策，再定位权威数据和代码；生成结果后检查文件存在性、哈希、程序输出和视觉结果。\\
2 & 以用户逐轮图表反馈为约束的修订提示 & 根据用户对建模示意图、数据密度、模型分面和径向网格数量的明确意见，仅修改图表脚本和图件，不改变题目数据口径。\\
3 & 以“先核验、再写作”的审计提示 & 对图4、图5、附录图 S2/S3 的源数据、参数、固定时间步和图件哈希做交叉检查；对模型假设中的潜热与温度处理回看 formal\_compute.py。\\
\bottomrule\end{longtable}}
\section*{表4 AI输出的采纳、人工修改和核验情况}
{\small\setlength{\tabcolsep}{3pt}\begin{longtable}{>{\raggedright\arraybackslash}p{0.06\linewidth}>{\raggedright\arraybackslash}p{0.29\linewidth}>{\raggedright\arraybackslash}p{0.15\linewidth}>{\raggedright\arraybackslash}p{0.4\linewidth}}
\toprule
序号 & AI生成内容概述 & 是否采纳 & 人工修改与核验情况\\
\midrule\endhead
1 & 数值计算和图表脚本草稿、调试建议 & 部分采纳 & 队员/主控逐项核对参数、边界条件、模型路线、输出 CSV 和运行日志；不符合冻结决策的内容不纳入活动结果。\\
2 & 图4三模型分面、图5六组网格、附录图 S2/S3 五组网格的图表排版 & 采纳并修改 & 根据用户反馈多轮调整；使用 check\_figure.py --min-dpi 300 --strict、源数据快照、PNG/SVG 哈希和人工视觉检查复核。\\
3 & 论文结构、图注和局限性表述 & 待队伍最终复核 & AI 只整理已批准证据；核心建模分析和最终表述由队伍审查，尤其保留“无潜热耦合、M2/M3按烘房温度处理”的限制。\\
4 & AI 使用声明和详情表 & 待队伍最终复核 & 按官网2026年试行规定逐项检查工具、用途、过程、采纳与人工核验，不使用“未使用AI”声明。\\
\bottomrule\end{longtable}}
\section*{截图说明}
本支撑材料未提供交互平台截图；2026年试行规定要求的是 AI 工具使用详情 PDF，并未将截图列为必需项。

\section*{规则来源}
全国大学生数学建模竞赛人工智能工具使用规定（2026年试行），官网发布日期 2026-08-03：

{\raggedright\url{https://www.mcm.edu.cn/html_cn/node/fef94648f2836ab6cc81586f4c38512b.html}\par}\end{document}